% Study sheet: Volume II, Chapter 4, Data Storage.
% Numbers: \val{...} from data_storage_studysheet_numbers.tex, written by data_storage_studysheet_numbers.py.
\documentclass{studysheet}
\input{data_storage_studysheet_numbers}
\sheetvolume{2}
\sheetchapter{4}{Data Storage}
\sheetsources{Numbers from the mlsysim registry and the chapter's scenarios; all units decimal (GB = 10$^9$ bytes); assumed inputs are marked.}

\begin{document}
% ===================================================================== FRONT
\frontside
\begin{bigidea}
The Fleet Law has no term for waiting on data: it assumes the \hl{next batch is already in
accelerator memory}. Storage keeps it by staging data close enough, and keeps progress by saving
state that \hl{can be restored}.
\end{bigidea}

\begin{leftcol}
\sheetsec{Key ideas}
\begin{ideas}
\idea{The I/O wall.}{Faster compute against a fixed input path speeds training only up to the delivery
  rate. Prefetching hides variation, never a sustained shortfall.}
\idea{Size from demand, then the slowest stage.}{Not from tier ratios: batch bytes per compute time,
  at every stage. GPUDirect Storage helps only if the host copy limits.}
\idea{Small files are a metadata problem.}{Signature: metadata operations at their limit while data
  servers idle. Packing thousands of samples per shard divides the opens.}
\idea{Commit boundaries make state safe.}{Pause for the local write; publish the manifest last; keep
  the previous generation until then. Replays write under stable keys.}
\idea{Total delivery cost decides placement.}{Capacity is billed per byte stored, egress per byte read,
  requests per operation. Data read often belongs close.}
\idea{Serving reads differently.}{A rollout is a burst of whole-model reads from every replica; graph
  retrieval is bound by random reads within a deadline.}
\end{ideas}

\vspace{16pt}
\sheetsec{Bandwidth by tier vs.\ what a node needs}
{\centering
\begin{tikzpicture}[sheet]
  \pgfmathsetmacro\k{17}
  \foreach \e in {-4,...,4}{\draw[gridline] ({100+(\e+4)*\k},-8) -- ({100+(\e+4)*\k},74);}
  \foreach \e/\t in {-3/{1 MB/s}, 0/{1 GB/s}, 3/{1 TB/s}}{\node[small, anchor=north] at ({100+(\e+4)*\k},-10) {\t};}
  \foreach \lab/\v/\y/\sty/\txt/\side/\dx in {%
      {HBM, per accelerator}/\val{ladderhbm}/68/white/{\val{hbmbandwidth} TB/s}/east/-4,
      {NVMe stripe, per node}/\val{ladderstripe}/52/white/{\val{stripebandwidth} GB/s}/west/4,
      {file-system share, per node}/\val{ladderpfs}/36/white/{\val{pfsshare} GB/s}/west/4,
      {image job needs, per node}/\val{ladderimage}/16/accent/{\val{imagenode} GB/s}/west/4,
      {text job needs, per node}/\val{laddertext}/0/accent/{\val{textnode} MB/s}/west/4}{
    \pgfmathsetmacro\x{100+(log10(\v)+4)*\k}
    \node[lab, anchor=east] at (94,\y) {\lab};
    \draw[dim] (100,\y) -- (236,\y);
    \filldraw[fill=\sty, draw=accent, line width=0.6pt] (\x,\y) circle (2.4pt);
    \node[note, fill=white, inner sep=1.2pt, anchor=\side] at (\x+\dx,\y) {\txt};}
  \draw[dim, dashed] (100,26) -- (236,26);
\end{tikzpicture}\par}
\vspace{3pt}
{\fontsize{8.4}{10.4}\selectfont\color{muted}Log scale; the chapter's jobs on \val{clusternodes} nodes sharing a
\val{pfsaggregate} TB/s file system. Demand, not HBM, sets the feed.\par}
\end{leftcol}%
\begin{rightcol}
\sheetsec{Formulas you will use}
\fcard{Input demand and the service chain}%
  {\begin{gathered}\RequiredInputBandwidth = \InputConsumerAcceleratorCount \, \TargetComputeTimeFraction \,
   \frac{\BatchDataVolume}{\IronLawComputeTime}\\[2pt]
   \InputPipelineServiceRate = \min_{\InputPipelineStageIndex} \InputStageServiceRate_{\InputPipelineStageIndex} \geq \RequiredInputBandwidth\end{gathered}}%
  {Accelerators reading distinct batches $\times$ target busy fraction $\times$ batch bytes per compute
   time. Every stage $\InputPipelineStageIndex$ must sustain it.}%
  {Compare stages in output-equivalent bytes.}{}
\fcard{Tier as a cache}%
  {\mathbb{E}[\TieredRequestServiceTime] = \HotTierHitProbability\,\HotTierServiceTime
   + (1-\HotTierHitProbability)\,\MissPathServiceTime}%
  {Hit share $\HotTierHitProbability$ at hot time $\HotTierServiceTime$; misses take $\MissPathServiceTime$.
   Delivered rate $\propto 1/\mathbb{E}[\TieredRequestServiceTime]$.}%
  {\val{chhit}\% hits, misses \val{chmiss}$\times$ slower: $\val{chmean}\,\HotTierServiceTime$, about
   1/\val{chfrac} of the all-hot rate.}{}
\fcard{Prefetch depth}%
  {\MinimumPrefetchBufferCount = \left\lceil \frac{\TailStepIoTime}{\IronLawComputeTime} \right\rceil}%
  {Ready batches that cover a p99 read delay. Host memory $=$ depth $\times$ batch $\times$ accelerators per node.}%
  {If reads are slower on average, an overlapped step takes the read time, the stall share is
   (read $-$ compute)$\,/\,$read, and no depth fixes it.}{}
\fcard{Checkpoint size and pause}%
  {\text{bytes} = \val{bytesperparam}\,\Parameters, \qquad
   \CheckpointWriteTime = \frac{\text{bytes} / \text{nodes}}{\text{local bandwidth}}}%
  {16-bit weights (2 bytes), 32-bit master weights (4), two Adam moments (4 each).}%
  {The pause is a barrier, a coordination cost in $\SynchronizationTime$; the background copy costs step
   time only if it contends with collectives or input reads.}{}
\end{rightcol}

% ===================================================================== BACK
\backside
\begin{answers}
\ans{Finish}{$\val{ckptgb} \div \val{sharebw}$ GB/s $=$ \val{directs} s, \val{directpct}\% of each
  \val{intervals} s interval, \val{directratio}$\times$ the staged \val{writes} s pause. Staged, the copy
  drains off the step; keep the previous generation until the manifest is published (checkpoint card).}
\ans{1}{$\val{qimages} \times \val{qkb}$ KB $=$ \val{qbatch} MB; $\val{qgpus} \times \val{qeta} \times
  \val{qbatch} \div$ \val{qms} ms $=$ \val{qneed} GB/s. Share $\val{qnodes} \times \val{pershare} \approx
  \val{qfair}$ GB/s suffices; decode, \val{qdecode} GB/s, is the minimum: not fed (demand card).}
\ans{2}{$\val{hit} + \val{missshare} \times \val{missx} = \val{meanservice}\,\HotTierServiceTime$: about
  \val{allhotpct}\% of the all-hot rate. Half the rate needs $\mathbb{E}[\TieredRequestServiceTime]$ of two
  hot times: \val{halfhit}\% hits (cache card).}
\columnbreak
\ans{3}{$\lceil \val{pnine} / \val{pms} \rceil = \val{depth}$ ready batches; $\val{depth} \times \val{pmb}
  \times \val{pgpus} = \val{pbuffer}$ MB per node. At \val{slowms} ms, stall \val{slowgap} $\div$ \val{slowms} $=$ \val{stallpct}\% at any depth (prefetch card).}
\ans{4}{Metadata: operations at their limit while data servers idle, and the lost time is
  input-queue wait, not collectives. Pack into shards: \val{filesper}$\times$ fewer opens (small files idea).}
\end{answers}

\begin{leftcol}
\sheetsec{Worked example}
\begin{practicebox}\boxtitle{How long does a 175B training job pause to save?}
\val{params}B parameters, \val{clusternodes} nodes, each with \val{drivecount} NVMe drives at
\val{drivebandwidth} GB/s; the shared file system absorbs \val{durablebw} TB/s; a save every
\val{intervalmin} minutes.\par\vspace{5pt}
{\centering
\begin{tikzpicture}[sheet]
  \pgfmathsetmacro\u{56}
  \pgfmathsetmacro\a{70+\val{writes}*\u}
  \pgfmathsetmacro\b{\a+\val{drains}*\u}
  \node[lab, anchor=east] at (66,25) {training};
  \node[lab, anchor=east] at (66,7) {shared copy};
  \draw[nbox] (70,20) rectangle (234,30);
  \draw[hotbox] (70,20) rectangle (\a,30);
  \node[note, anchor=south west] at (70,31) {pause \val{writes} s};
  \draw[accbox] (\a,2) rectangle (\b,12);
  \node[lab] at ({(\a+\b)/2},7) {background, \val{drains} s};
  \draw[accent, line width=1pt] (\b,-1) -- (\b,15);
  \node[note, anchor=north east] at (\b+1,-2) {manifest last};
\end{tikzpicture}\par}\vspace{3pt}
\begin{steps}
\item Size: \val{params}B $\times$ \val{bytesperparam} bytes $=$ \val{ckptgb} GB, written once in total.
\item Per node: \val{ckptgb} $\div$ \val{clusternodes} $=$ \val{shardgb} GB at \val{localbw} GB/s
  $=$ \res{\val{writes} s pause} (\val{pausepct}\% of the interval).
\item Background: \val{ckptgb} GB $\div$ \val{durablebw} TB/s $=$ \res{\val{drains} s} to drain, off the step.
\end{steps}
\takeaway{The pause ends at the local write; durability comes later.}
\end{practicebox}

\vspace{12pt}
\sheetsec{Your turn: finish it}
\begin{practicebox}\boxtitle{What if the save goes straight to a busy file system?}
Same save, but \val{jobs} jobs share the file system, so this one gets
\val{durablebw} TB/s $\div$ \val{jobs} $=$ \val{sharebw} GB/s.\par\vspace{4pt}
\textbf{Given:} the checkpoint is \val{ckptgb} GB.\par
\textbf{You:} find the pause when every byte goes directly to the file system, its share of each
\val{intervalmin}-minute interval, and what local staging changes.
\end{practicebox}
\end{leftcol}%
\begin{rightcol}
\sheetsec{Common mistakes}
\mistake{``HBM is \val{tierratio}$\times$ faster than an NVMe drive, so storage can never feed the accelerators.''}%
  {That ratio compares tiers, not demand; HBM mostly carries reused weights and activations.
   The chapter's text job needs \val{textjob} MB/s for \val{textgpus} GPUs.}{}
\mistake{``Object storage is the cheapest tier, so stream from it every epoch.''}%
  {Egress is billed per byte read: a \val{corpustb} TB corpus read \val{epochs} times at \$\val{egress}/GB
   costs \$\val{perrun} per run, against \$\val{storeyear} a year to store it. Staging on local NVMe reads it
   across the boundary once per run: \$\val{stageegress}.}{}
\mistake{``If a pipeline worker crashes, replay its work; one extra write does no harm.''}%
  {An appended replay writes the row twice, so the event counts twice in every training set. Write
   under a stable key (data version, event ID) so a replay overwrites the same row.}{}

\vspace{2pt}
\sheetsec{Check yourself}
\begin{qs}
\item \assumed{} \val{qgpus} GPUs each read \val{qimages} images of \val{qkb} KB per \val{qms} ms step;
  target \val{qeta}. What bandwidth is needed? Its \val{qnodes} nodes get the per-node file-system share;
  decoding sustains \val{qdecode} GB/s. Is it fed?
\item \assumed{} \val{hitpct}\% of reads hit local NVMe and misses take \val{missx}$\times$ as long. What
  share of the all-hot rate is delivered? What hit rate gives half?
\item \assumed{} compute takes \val{pms} ms, file-system p99 reads \val{pnine} ms, decoded batches
  \val{pmb} MB, \val{pgpus} GPUs per node. Depth and host memory? If reads average \val{slowms} ms, does a deeper queue help?
\item A job moves from \val{shards} shards to \val{files} million single-image files on a \val{ratedbw} GB/s file
  system and gets under one percent of it. Bandwidth or metadata? What evidence decides? The fix?
\end{qs}
\end{rightcol}
\end{document}
